% 07-memory.tex
\section{内存：DataPtr、Storage、TensorImpl 和 Tensor}
\label{sec:memory}

TensorPlay 中的内存是严格的三层堆栈：原始分配（\texttt{DataPtr} 加上每个设备 \texttt{Allocator}）、字节容器（\texttt{StorageImpl} / \texttt{Storage}）以及具有形状和调度状态的类型化视图（\texttt{TensorImpl} / \texttt{Tensor}）。每层都是物理隔离的，在下面的 \texttt{path:line} 中引用，并且设计成最便宜的操作（视图）根本不接触分配器。

\subsection{三层和所有权}
\label{sec:memlayers}

\begin{figure}[H]
\centering
\begin{tikzpicture}[every node/.style={font=\scriptsize}, node distance=0.45cm]
\node[libbox, fill=blue!8, draw=tpblue, minimum width=3.4cm, minimum height=0.9cm] (dptr) {\texttt{DataPtr} \\};
\node[libbox, fill=orange!8, draw=tporange, below=of dptr, minimum width=3.4cm, minimum height=0.9cm] (alloc) {\texttt{Allocator} / \texttt{CachingAllocator} \\};
\node[libbox, fill=green!8, draw=tpgreen, below=of alloc, minimum width=5.8cm, minimum height=1.05cm] (simpl) {\texttt{StorageImpl} \{ \texttt{data\_ptr}, \texttt{nbytes}, \texttt{allocator*}, \texttt{resizable} \} \\ \quad $\rightarrow$ \quad \texttt{Storage} \{ \texttt{shared\_ptr<StorageImpl>} \} \\};
\node[libbox, fill=purple!8, draw=tppurple, below=of simpl, minimum width=11cm, minimum height=1.7cm] (timpl) {\texttt{TensorImpl} \{ 8 groups: offset / shape+strides / dtype+device / version+inference / contig+format / autograd+view / storage+shared+sparse / transform+quantizer \} \\ \quad --- \texttt{SizesAndStrides} inline \texttt{[5]} \quad --- \texttt{VariableVersion} \texttt{atomic<uint32\_t>}};
\node[libbox, fill=gray!10, draw=gray, below=of timpl, minimum width=6.2cm, minimum height=0.9cm] (tensor) {\texttt{Tensor} \{ \texttt{shared\_ptr<TensorImpl> impl\_} \} + 479 generated methods \\ \texttt{Tensor.h}};
\draw[arrow] (dptr) -- (alloc);
\draw[arrow] (alloc) -- (simpl);
\draw[arrow] (simpl) -- (timpl);
\draw[arrow] (timpl) -- (tensor);
\node[right=0.35cm of simpl, align=left, font=\tiny, text=gray] (note1) {\texttt{Storage::is\_same} = pointer \\\texttt{identity on StorageImpl*}};
\node[right=0.35cm of timpl, align=left, font=\tiny, text=gray] (note2) { aliases\\\texttt{Storage}, shares counter\\ copies bytes};
\draw[dashed, color=gray] (note1) -- (simpl);
\draw[dashed, color=gray] (note2) -- (timpl);
\end{tikzpicture}
\caption{三层：\texttt{DataPtr}（原始，仅移动）$\rightarrow$ \texttt{StorageImpl/Storage}（字节，引用计数）$\rightarrow$ \texttt{TensorImpl/Tensor}（类型化，成形，版本化）。视图是零拷贝别名；克隆分配底层。}
\label{fig:memlayers}
\end{figure}

只有 \texttt{Tensor} 可供用户复制； \texttt{DataPtr} 只能移动，\texttt{DataPtr} 只能移动，而 \texttt{StorageImpl} 永远不会通过值暴露。 \texttt{TensorImpl} 由 \texttt{shared\_ptr} 管理，因此视图、稀疏组件和批处理包装器可以为其别名，而无需手动生命周期跟踪。

\subsection{DataPtr：带有已擦除删除器和设备的仅移动句柄}
\label{sec:dataptr}

将 \texttt{DataPtr} 定义为四个字上的仅移动句柄：

\begin{lstlisting}[caption={move-only handle, 4 words.}, label=lst:dataptrfull]
class DataPtr {
 private:
  void* data_;          // typed data pointer returned to callers
  void* ctx_;           // context passed to deleter (not necessarily data_)
  DeleterFnPtr deleter_; // void(*)(void*)  --- capture-less, trivially copyable
 public:
  Device device_;        // Device{type,index} travels with the allocation
  DataPtr() : data_(nullptr), ctx_(nullptr), deleter_(nullptr),
              device_(DeviceType::Unknown) {}
  DataPtr(void* data, void* ctx, DeleterFnPtr ctx_deleter, Device device)
      : data_(data), ctx_(ctx), deleter_(ctx_deleter), device_(device) {}
  DataPtr(void* data, DeleterFnPtr deleter, Device device)  // ctx == data
      : data_(data), ctx_(data), deleter_(deleter), device_(device) {}
  DataPtr(DataPtr&& other) noexcept : DataPtr() { swap(other); }
  DataPtr& operator=(DataPtr&& other) noexcept { if(this!=&other){clear(); swap(other);} return *this; }
  DataPtr(const DataPtr&) = delete;             // no copy
  DataPtr& operator=(const DataPtr&) = delete;
  ~DataPtr(){ clear(); }
  void clear(){ if(deleter_) deleter_(ctx_); data_=ctx_=nullptr; deleter_=nullptr; }
  void swap(DataPtr& other) noexcept;           // swaps 4 words
  void* get(){ check_device(); return data_; }
  const void* get() const { check_device(); return data_; }
  void* get_context() const { return ctx_; }
  DeleterFnPtr get_deleter() const { return deleter_; }
  void* release(){ void* p=data_; data_=ctx_=nullptr; deleter_=nullptr; return p; }
  operator bool() const { return data_!=nullptr; }
};
\end{lstlisting}

设计点，每个点都可以在以下方面验证：

\begin{itemize}[leftmargin=1.2em]
\item \textbf{Erased deleter via ->.} 别名 \texttt{DeleterFnPtr = void(*)(void*)} 是一个普通函数指针，而不是 \texttt{std::function}。有状态清理是通过将单独的 \texttt{ctx\_} 传递到 \texttt{deleter\_(ctx\_)} 来编码的——这是一个避免分配和捕获的上下文承载签名。替代构造函数 \texttt{DataPtr(void*,DeleterFnPtr,Device)} 为无状态情况（例如 \texttt{deleteCPP}）设置 \texttt{ctx\_==data\_}。
\item \textbf{Device travels with the allocation.} 公共字段 \texttt{Device device\_} 被初始化为 \texttt{Unknown} 以表示空状态，并与指针一起交换。当 \texttt{NDEBUG} 不存在时，每个 \texttt{get} 都会强制进行调试检查，并对 \texttt{Unknown} 访问发出警告。较高层（\texttt{StorageImpl}、\texttt{TensorImpl}）读取 \texttt{data\_ptr.device\_} 来决定复制路径，而不是缓存该类型的第二个副本。
\item \textbf{Move-only, not copyable.} 两个复制操作都是 \texttt{=delete}，并且两个移动操作都通过空默认值进行交换。表~\ref{tab:dataptr} 总结了标题中提供的删除器选择。
\end{itemize}

\begin{table}[H]
\centering
\caption{---删除器品种及用途。全部都是 \texttt{DeleterFnPtr}；州政府乘坐 \texttt{ctx\_}。}
\label{tab:dataptr}
\small
\begin{tabular}{lll}
\toprule
象征 & 签名/地点 & 使用时 \\
\midrule
\texttt{deleteNothing} & \texttt{void(void*)} & 借用外部存储器（无所有权） \\
\texttt{deleteCPP} & \texttt{delete[] char*} & CPU 测试中的后备路径 \\
\texttt{CachingAllocator::deleter} & \texttt{void(void*)} & 正常路径：\texttt{ctx} 是数据指针，标头位于 \texttt{ptr-64} \\
\texttt{release} & \texttt{void*} & 调用者取得所有权，在不调用它的情况下删除删除器 \\
\bottomrule
\end{tabular}
\end{table}

大小含义是直接的：\texttt{sizeof(DataPtr)} 是四个指针（三个 \texttt{void*}/函数指针加上 \texttt{Device} 带填充），可以通过 \texttt{swap} 轻松重定位，并且可以安全地从 \texttt{Allocator::allocate} 按值返回。

\subsection{分配器：带有 64 字节标头、\texttt{free\_blocks\_} 和副本调度的分桶缓存}
\label{sec:allocator}

\subsubsection{界面}

是抽象接口：

\begin{lstlisting}[caption={abstract surface.}, label=lst:allociface]
class P10_API Allocator {
 public:
  virtual ~Allocator() = default;
  virtual DataPtr allocate(size_t nbytes) const = 0;
  virtual DataPtr allocate(size_t nbytes, const Device& device) const {
    (void)device; return allocate(nbytes);
  }
};
P10_API Allocator* getAllocator(DeviceType t);   // dispatch by type
P10_API Allocator* getCPUAllocator();             // singleton accessor
P10_API void copyAllocationBytes(void* dst, const Device& dst_dev,
                                 const void* src, const Device& src_dev,
                                 size_t nbytes);
\end{lstlisting}

\texttt{getCPUAllocator}、\texttt{getAllocator(DeviceType)} 和 \texttt{copyAllocationBytes} 完成表面。特定于设备的构建公开了其他符号，例如 \texttt{getPinnedMemoryAllocator} 和 \texttt{getCUDAAllocator}。

\subsubsection{缓存分配器内部结构}

CPU实现完全存在于一个匿名命名空间中。其结构：

\begin{lstlisting}[caption={bucketed caching allocator.}, label=lst:cachingalloc]
class CachingAllocator : public Allocator {
  struct Header { size_t size; };                // raw allocation size incl. header
  static constexpr size_t HEADER_SIZE = 64;      // 64-byte alignment
  mutable std::mutex mutex_;
  mutable std::unordered_map<size_t, std::vector<void*>> free_blocks_; // bucket -> raw ptrs
 public:
  static CachingAllocator* instance(){ static CachingAllocator* inst=new CachingAllocator(); return inst; }
  static void deleter(void* ptr){ // ptr == data pointer (header+64)
    char* raw = static_cast<char*>(ptr) - HEADER_SIZE;
    Header* h = reinterpret_cast<Header*>(raw);
    size_t size = h->size;                       // rounded size, incl. header
    prof::mem_record_free(ptr, int64_t(size-HEADER_SIZE), false, -1, -1);
    instance()->free(raw, size);
  }
  void free(void* ptr, size_t size){ std::lock_guard<std::mutex> l(mutex_); free_blocks_[size].push_back(ptr); }
  DataPtr allocate(size_t nbytes) const override {
    size_t total_size = nbytes + HEADER_SIZE;    // data starts at raw+64
    total_size = (total_size + 63) & ~size_t(63); // bucket to 64 bytes
    void* ptr=nullptr;
    { std::lock_guard<std::mutex> l(mutex_);
      auto it=free_blocks_.find(total_size);
      if(it!=free_blocks_.end() && !it->second.empty()){ ptr=it->second.back(); it->second.pop_back(); }
    }
    if(!ptr) ptr = alloc_aligned(total_size);    // posix_memalign 64
    if(!ptr) TP_THROW(RuntimeError,"Out of memory");
    reinterpret_cast<Header*>(ptr)->size = total_size;
    void* data_ptr = static_cast<char*>(ptr) + HEADER_SIZE;
    prof::mem_record_alloc(data_ptr, int64_t(total_size-HEADER_SIZE), false, -1, -1);
    return DataPtr(data_ptr, deleter, Device(DeviceType::CPU));
  }
};
\end{lstlisting}

具体细节：

\begin{itemize}[leftmargin=1.2em]
\item \textbf{64-byte header and alignment.} \texttt{HEADER\_SIZE=64} 保证每当 \texttt{raw} 时 \texttt{data = raw+64} 都是 64 字节对齐的，因为 64 划分了对齐方式。帮助程序 \texttt{alloc\_aligned} 在 POSIX 上使用 \texttt{posix\_memalign(\&ptr,64,nbytes)}，在 Windows 上使用 \texttt{\_aligned\_malloc}。
\item \textbf{Bucketing.} 总大小 \texttt{nbytes+64} 向上舍入为 64 的下一个倍数：\texttt{(total\_size+63) \& \~{}63}。存储桶通过此舍入总数 (\texttt{free\_blocks\_.find(total\_size)} at) 进行键控，减少碎片，同时保持重用检查 O(1)。 \texttt{Header::size} 字段存储舍入大小，因此 \texttt{deleter} 可以将块返回到正确的存储桶，而无需重新计算。
\item \textbf{Free list.} \texttt{free\_blocks\_: unordered\_map<size\_t, vector<void*>>} 是每个桶的堆栈； \texttt{free} 压入原始指针，\texttt{allocate} 在存在时弹出它。地图和向量在两条路径上均由 \texttt{mutex\_} 保护。
\item \textbf{Deleter path.} \texttt{deleter(void* ptr)} 通过 \texttt{ptr-64} 恢复标头，读取 \texttt{size}，为探查器记录空闲，然后重新插入到缓存中，而不是调用 \texttt{free}。因此 \texttt{allocate} 返回的 \texttt{DataPtr} 带有 \texttt{deleter==CachingAllocator::deleter} 和 \texttt{ctx==data}，因此销毁会回收而不是释放到操作系统。
\item \textbf{Profiler hook.} 分配和释放都使用桶舍入的用户大小 \texttt{total\_size-64} 调用 \texttt{free}，保持捕获的分配/释放字节一致。
\end{itemize}

\begin{table}[H]
\centering
\caption{\texttt{CachingAllocator} 常量和状态。全部私人；通过 \texttt{instance} 单例。}
\label{tab:cachingconst}
\small
\begin{tabular}{lll}
\toprule
成员 & 定义 & 值/类型 \\
\midrule
\texttt{HEADER\_SIZE} &  & \texttt{64} 字节（标头大小和对齐方式） \\
\texttt{Header::size} &  & 包括标题在内的四舍五入总计 (\texttt{nbytes+64} $\uparrow$ 64) \\
\texttt{free\_blocks\_} &  & \texttt{unordered\_map<size\_t, vector<void*>>}：桶 $\rightarrow$ 原始块 \\
\texttt{mutex\_} &  & 保护 alloc 和 free 上的免费地图 \\
\texttt{alloc\_aligned} &  & \texttt{posix\_memalign(64)} / \texttt{\_aligned\_malloc} \\
\texttt{deleter} &  & 恢复 \texttt{ptr-64} 处的标头，推回正确的存储桶 \\
\bottomrule
\end{tabular}
\end{table}

\subsubsection{全球调度：\texttt{getCPUAllocator}、\texttt{getAllocator}、\texttt{copyAllocationBytes}}

\texttt{getCPUAllocator} 返回 \texttt{CachingAllocator::instance}，一个泄漏的单例，因此它的析构函数永远不会与后期 \texttt{DataPtr} 销毁竞争。 \texttt{getAllocator(DeviceType)} 打开类型并为未构建的后端抛出 \texttt{NotImplementedError}，并为 \texttt{USE\_CUDA} 和 \texttt{USE\_VULKAN} 进行条件编译。

\texttt{copyAllocationBytes} 是 \texttt{StorageImpl::set\_nbytes} 使用的跨设备 \texttt{memcpy}：

\begin{lstlisting}[caption={copy dispatch.}, label=lst:copybytes]
void copyAllocationBytes(void* dst, const Device& dst_dev,
                         const void* src, const Device& src_dev, size_t nbytes){
  if(!dst||!src||nbytes==0) return;
  if(dst_dev.is_cpu() && src_dev.is_cpu()){ std::memcpy(dst,src,nbytes); return; }
#ifdef USE_VULKAN
  if(dst_dev.is_vulkan()||src_dev.is_vulkan()){
    vulkan::copyHostVisibleBytes(const_cast<void*>(dst),dst_dev,
                                 const_cast<void*>(src),src_dev,nbytes); return;
  }
#endif
#ifdef USE_CUDA
  const Device cuda_dev = dst_dev.is_cuda()? dst_dev : src_dev;
  cuda::CUDAGuard guard(cuda_dev);
  auto stream = cuda::getCurrentCUDAStream(int(cuda_dev.index()));
  cudaMemcpyKind kind = cudaMemcpyDefault;
  if(dst_dev.is_cuda() && src_dev.is_cuda()) kind=cudaMemcpyDeviceToDevice;
  else if(dst_dev.is_cuda()) kind=cudaMemcpyHostToDevice;
  else kind=cudaMemcpyDeviceToHost;
  cuda::checkCuda(cudaMemcpyAsync(dst,src,nbytes,kind,stream.stream()),"cudaMemcpyAsync (storage resize)");
  if(dst_dev.is_cuda()) cuda::recordStream(dst,stream);
  if(src_dev.is_cuda()) cuda::recordStream(const_cast<void*>(src),stream);
  if(dst_dev.is_cpu()||src_dev.is_cpu()) stream.synchronize();
  return;
#else
  if(dst_dev.is_cuda()||src_dev.is_cuda()) TP_THROW(RuntimeError,"cannot copy CUDA storage in CPU-only build");
#endif
  TP_THROW(RuntimeError,"cannot copy storage between these devices");
}
\end{lstlisting}

这三个分支是具体的：CPU--CPU 使用 \texttt{std::memcpy}，Vulkan 主机可见字节使用专用路径，CUDA 在当前流上使用 \texttt{cudaMemcpyAsync} 进行流记录，并且当任一侧是 CPU 时进行主机可见同步。这是 CPU 与加速器副本唯一存在分歧的地方；高层盲目调用。

\subsection{StorageImpl：字节加上 \texttt{resizable} 和 \texttt{Storage}：\texttt{shared\_ptr} 包装器}
\label{sec:storage}

\subsubsection{存储实现}

是字节容器：

\begin{lstlisting}[caption={bytes, allocator, resizable.}, label=lst:storageimplfull]
struct P10_API StorageImpl {
  DataPtr data_ptr;       // owns the allocation, carries Device
  size_t nbytes;          // logical byte extent (may be < bucket size)
  Allocator* allocator;   // non-owning; used for growth
  bool resizable;         // false for wrapped external memory
  StorageImpl(size_t size, Allocator* alloc, bool resizable=true) // :19
      : nbytes(size), allocator(alloc), resizable(resizable) {
    if(allocator) data_ptr = allocator->allocate(nbytes);
  }
  StorageImpl(size_t size, Allocator* alloc, const Device& device, bool resizable=true) // :26
      : nbytes(size), allocator(alloc), resizable(resizable) {
    if(allocator) data_ptr = allocator->allocate(nbytes, device);
  }
  StorageImpl(DataPtr&& ptr, size_t size, Allocator* alloc, bool resizable=false) // :33
      : data_ptr(std::move(ptr)), nbytes(size), allocator(alloc), resizable(resizable) {}
  void* data(){ return data_ptr.get(); }
  const void* data() const { return data_ptr.get(); }
  void set_nbytes(size_t new_nbytes){ // :44
    if(nbytes==new_nbytes) return;
    if(!resizable || !allocator) TP_THROW(RuntimeError,"Storage is not resizable");
    DataPtr new_data = allocator->allocate(new_nbytes, data_ptr.device_); // :50
    size_t copy_size = std::min(nbytes, new_nbytes);                        // :52
    if(data_ptr.get() && new_data.get())
      copyAllocationBytes(new_data.get(), new_data.device_,                 // :54
                          data_ptr.get(), data_ptr.device_, copy_size);
    data_ptr = std::move(new_data); nbytes = new_nbytes;                    // :58
  }
};
\end{lstlisting}

关键语义：

\begin{itemize}[leftmargin=1.2em]
\item \textbf{Four fields, no shape.} \texttt{data\_ptr}、\texttt{nbytes}、\texttt{allocator*}、\texttt{resizable} 是整个状态。解释为 \texttt{float32} 与\ \texttt{int64} 是 \texttt{TensorImpl::dtype\_} 的工作，而不是存储。
\item \textbf{Three constructors.} 通过分配器分配、在设备上分配和包装现有-\texttt{DataPtr}。包装形式默认为 \texttt{resizable=false} 因为外部拥有的删除器不能由分配器重新生成。两个分配构造函数在调用 \texttt{allocator->allocate(nbytes)} / \texttt{allocator->allocate(nbytes,device)} 之前都会保护 \texttt{if(allocator)}。
\item \textbf{Growth via ->.} \texttt{StorageImpl::set\_nbytes} 在同一设备上分配新大小的新 \texttt{DataPtr}，通过 \texttt{copyAllocationBytes} 复制 \texttt{min(old,new)} 字节，然后交换。当 \texttt{!resizable || !allocator} 时它会抛出。无操作路径 \texttt{nbytes==new\_nbytes} 提前返回而不进行分配。
\end{itemize}

\subsubsection{储物手柄}

是一个 \texttt{shared\_ptr} 包装器：

\begin{lstlisting}[caption={ref-counted handle.}, label=lst:storagehandle]
class P10_API Storage {
 public:
  Storage() = default;
  explicit Storage(size_t nbytes, Allocator* allocator=nullptr){ // :14
    if(!allocator) allocator=getCPUAllocator();
    impl_=std::make_shared<StorageImpl>(nbytes, allocator);
  }
  Storage(size_t nbytes, Allocator* allocator, const Device& device){ // :19
    if(!allocator) allocator=getCPUAllocator();
    impl_=std::make_shared<StorageImpl>(nbytes, allocator, device);
  }
  Storage(DataPtr&& data_ptr, size_t nbytes, Allocator* allocator=nullptr){ // :25
    impl_=std::make_shared<StorageImpl>(std::move(data_ptr), nbytes, allocator, false);
  }
  void* data() const { return impl_? impl_->data() : nullptr; }
  template<typename T> T* data() const { return static_cast<T*>(data()); }
  size_t nbytes() const { return impl_? impl_->nbytes : 0; }
  Device device() const { return impl_? impl_->data_ptr.device_ : Device(DeviceType::CPU); }
  bool defined() const { return impl_!=nullptr; }
  bool is_same(const Storage& other) const { return impl_==other.impl_; } // :47 identity
  bool resizable() const { return impl_ && impl_->resizable; }
  void set_nbytes(size_t nn){ if(!impl_) impl_=std::make_shared<StorageImpl>(nn, getCPUAllocator()); else impl_->set_nbytes(nn); }
  Allocator* allocator() const { return impl_? impl_->allocator : nullptr; }
  long use_count() const { return impl_.use_count(); }
  const std::shared_ptr<StorageImpl>& unsafeGetStorageImpl() const { return impl_; } // :67 backend only
 private:
  std::shared_ptr<StorageImpl> impl_; // :76
};
\end{lstlisting}

对于视图和克隆来说重要的后果：

\begin{itemize}[leftmargin=1.2em]
\item \textbf{View identity is pointer identity.} \texttt{is\_same} 比较 \texttt{shared\_ptr} 地址，而不是字节内容。 \texttt{b=a.view([3,2])} 的 Python 测试 \texttt{a.untyped\_storage.data\_ptr==b.untyped\_storage.data\_ptr} 是正确的，因为两个 \texttt{TensorImpl} 对象共享相同的 \texttt{Storage::impl\_}。
\item \textbf{Default allocator.} 每个 \texttt{Storage(size\_t, Allocator*)} 构造函数在传递 null 时都会回退到 \texttt{getCPUAllocator}，因此省略分配器的调用者仍然尊重存储桶缓存。
\item \textbf{Wrap is non-resizable.} \texttt{Storage(DataPtr\&\&,...)} 将 \texttt{resizable=false} 转发给 \texttt{StorageImpl}，匹配 \texttt{DataPtr(DataPtr\&\&)} 的外部所有权合约。
\end{itemize}

\subsection{TensorImpl：八组、386 行、一个扁平结构}
\label{sec:tensorimpl}

是一个具有 \texttt{Storage}、shape、dtype、设备、版本控制、布局、autograd、视图标识、共享 oneDNN 状态、稀疏、vmap 和量化状态的单个类。这些字段分为八个逻辑组，直接映射到成员声明：

\begin{table}[H]
\centering
\caption{---八组。订单与声明相符；明确显示默认值。}
\label{tab:tensorimplfields}
\small
\begin{tabular}{p{2.6cm}p{4.7cm}p{6.2cm}}
\toprule
团体 & 成员（线） & 语义学 \\
\midrule
1. 偏移量 & \texttt{storage\_offset\_:size\_t} & 元素在 \texttt{storage\_} 中的偏移量（以元素为单位，而不是以字节为单位）。按 \texttt{itemsize} 英寸缩放。 \\
2. 形状+步幅 & \texttt{sizes\_and\_strides\_(SizesAndStrides)} & 逻辑尺寸和步幅；内联排名 $\le 5$ （请参阅 \S\ref{sec:sizesstrides}）。 \\
3.标量类型+设备 & \texttt{dtype\_:DType}, \texttt{device\_:Device} & 元素类型和放置。 \texttt{itemsize} 是 \texttt{elementSize(dtype\_)}。 \\
4.版本+推理 & \texttt{version\_counter\_:VariableVersion}, \texttt{inference\_tensor\_:bool} & 原子版本可通过视图别名；推理标志禁用碰撞。 \\
5. 布局 & \texttt{is\_contiguous\_:bool}, \texttt{memory\_format\_:MemoryFormat} & 行主 vs.\ 通道最后标签； \texttt{is\_contiguous\_in(format)} at 将步幅与规范进行比较。 \\
6. Autograd + 查看标志 & \texttt{autograd\_meta\_:shared\_ptr<AutogradMetaBase>}, \texttt{is\_view\_:bool} & 每个张量的梯度状态（从未复制，请参阅 \S\ref{sec:copyctor}）；查看 \texttt{detach\_} 检查位。 \\
7、存储层 & \texttt{storage\_:Storage}, \texttt{SharedState \{onednn\_md, onednn\_memory\_cache\}}, \texttt{SparseState \{layout, indices, values, crow, col, sparse\_sizes, coalesced\}} & \texttt{storage\_} 是引用计数的字节； \texttt{SharedState} 缓存 oneDNN 描述符； \texttt{SparseState} 拥有 COO/CSR 组件。 \\
8. 变换+量化 & \texttt{transform\_value\_:shared\_ptr<TensorImpl>}, \texttt{transform\_batch\_dim\_, transform\_level\_}, \texttt{quantizer\_:shared\_ptr<Quantizer>} & Vmap 包装器（物理 vs.\ 逻辑）和仿射量化器（共享、不可变）。 \\
\bottomrule
\end{tabular}
\end{table}

类头的浓缩视图：

\begin{lstlisting}[caption={eight groups in one class.}, label=lst:tensorimplfull]
class P10_API TensorImpl {
 private:
  size_t storage_offset_;                          // :30 group 1
  SizesAndStrides sizes_and_strides_;              // :31 group 2
  DType dtype_; Device device_;                    // :32-33 group 3
  VariableVersion version_counter_{!InferenceMode::is_enabled()}; // :34 group 4
  bool inference_tensor_ = InferenceMode::is_enabled();           // :35
  bool is_contiguous_; MemoryFormat memory_format_ = Contiguous;  // :37,41 group 5
  std::shared_ptr<AutogradMetaBase> autograd_meta_; bool is_view_=false; // :45,49 group 6
  Storage storage_;                                              // :54 group 7
  struct SharedState{ std::shared_ptr<void> onednn_md; std::shared_ptr<void> onednn_memory_cache; } shared_state_; // :56
  struct SparseState{ int layout=0; std::shared_ptr<TensorImpl> indices,values,crow,col; std::vector<int64_t> sparse_sizes; bool coalesced=false; } sparse_state_; // :68
  std::shared_ptr<TensorImpl> transform_value_; int64_t transform_batch_dim_=-1, transform_level_=-1; // :82 group 8
  std::shared_ptr<Quantizer> quantizer_;                                         // :89
 public:
  // ctors, accessors, version, layout, sparse, batch, quant --- see listings below
};
\end{lstlisting}

选定的访问器说明了分层：

\begin{itemize}[leftmargin=1.2em]
\item \texttt{storage} / \texttt{set\_storage} / \texttt{storage\_offset} 公开 \texttt{Storage} 句柄和元素偏移量。将 \texttt{storage\_offset\_ * itemsize} 添加到基字节指针，并且在 CUDA 上，在返回之前调用 \texttt{cuda::recordStream(base\_ptr, device\_)}。
\item \texttt{sizes\_and\_strides}、\texttt{sizes}、\texttt{strides}、\texttt{dim}、\texttt{numel} 转发到 \texttt{SizesAndStrides}。动态计算调度键，在批处理时单独返回 vmap 键（请参阅 \S\ref{sec:vmap}）。
\item \texttt{is\_contiguous}、\texttt{memory\_format}、\texttt{set\_memory\_format}、\texttt{is\_contiguous\_in(format)} 维护布局标签。
\item \texttt{share\_version\_counter}、\texttt{bump\_version}、\texttt{reset\_version} 将别名连接到突变跟踪。
\end{itemize}

\subsection{SizesAndStrides：内联排名 $\le 5$}
\label{sec:sizesstrides}

是形状容器：

\begin{lstlisting}[caption={inline up to 5 dims.}, label=lst:sizesstrides]
class P10_API SizesAndStrides {
 public:
  static constexpr size_t kInlineSize = 5;             // :18
  SizesAndStrides() = default;                          // 0-dim scalar
  explicit SizesAndStrides(const std::vector<int64_t>& sizes);           // contiguous strides
  SizesAndStrides(const std::vector<int64_t>& sizes, const std::vector<int64_t>& strides);
 private:
  bool is_heap() const { return size_ > kInlineSize; } // :117
  const int64_t* sizes_data() const { return is_heap()? heap_sizes_.get() : inline_sizes_; } // :118
  size_t size_ = 0;                                     // :127 rank
  int64_t inline_sizes_[kInlineSize] = {};              // :128
  int64_t inline_strides_[kInlineSize] = {};            // :129
  std::unique_ptr<int64_t[]> heap_sizes_;   // :130 null while rank <=5
  std::unique_ptr<int64_t[]> heap_strides_; // :131
};
\end{lstlisting}

实施于：

\begin{lstlisting}[caption={extent transition.}, label=lst:setextent]
void SizesAndStrides::set_extent(size_t new_size){
  const bool was_heap = is_heap(); const bool will_be_heap = new_size > kInlineSize; // :75-76
  if(will_be_heap){ if(!was_heap || size_!=new_size){ heap_sizes_.reset(new int64_t[new_size]); heap_strides_.reset(new int64_t[new_size]); } }
  else if(was_heap){ heap_sizes_.reset(); heap_strides_.reset(); } // :82 frees heap when shrinking
  size_ = new_size;
}
\end{lstlisting}

结果：

\begin{itemize}[leftmargin=1.2em]
\item 大多数张量（等级 $\le 5$）从不分配：大小和步幅都位于两个 \texttt{[5]} 内联数组中。在这种情况下，\texttt{heap\_sizes\_}/\texttt{heap\_strides\_} 对为空，\texttt{sizes\_data}/\texttt{strides\_data} 在 \texttt{is\_heap} 上分支。
\item 复制和移动使用 \texttt{swap} ，以便移出的堆张量干净地转移其数组的所有权。复制构造函数根据 \texttt{is\_heap} 复制内联数组或堆分配。
\item 连续的步长是 \texttt{compute\_contiguous\_strides}：具有 \texttt{strides.back=1} 和 \texttt{strides[i]=strides[i+1]*max(size[i+1],1)} 的行优先，因此形状 \texttt{(2,0)} 得到 \texttt{(1,1)}。视图步幅使用 \texttt{compute\_view\_strides} ，它将旧形状分割成连续的块，并要求新形状平铺相同的块；不兼容的视图返回 \texttt{std::nullopt}。 \texttt{clone\_impl} 使用非重叠密集检查来决定是否保留步幅。
\end{itemize}

\begin{table}[H]
\centering
\caption{\texttt{SizesAndStrides} --- 内联与\ 堆。}
\label{tab:sizesinline}
\small
\begin{tabular}{lll}
\toprule
财产 & 内联路径（等级 $\le 5$） & 堆路径（等级 $>5$） \\
\midrule
贮存 & \texttt{inline\_sizes\_[5]}, \texttt{inline\_strides\_[5]} & \texttt{heap\_sizes\_}, \texttt{heap\_strides\_} \\
ctor/copy 上的分配 & 0 & 2 $\times$ \texttt{new int64\_t[rank]} \\
谓词 & \texttt{!is\_heap} & \texttt{is\_heap} \\
过渡 & \texttt{set\_extent} 在收缩到 $\le 5$ 时释放堆 & \texttt{set\_extent} 在增长超过 5 时分配堆 \\
\bottomrule
\end{tabular}
\end{table}

\subsection{VariableVersion：Eager Atomic Counter，由视图共享}
\label{sec:version}

是突变追踪器：

\begin{lstlisting}[caption={eager atomic counter.}, label=lst:varversion]
struct P10_API VersionCounter {            // :13
  std::atomic<uint32_t> version_{0};      // :14
  void bump() noexcept { version_.fetch_add(1, std::memory_order_relaxed); } // :16
  uint32_t current_version() const noexcept { return version_.load(std::memory_order_relaxed); } // :17
};
class P10_API VariableVersion {            // :24 handle
  std::shared_ptr<VersionCounter> counter_; // :26 shared among aliases
  bool enabled_ = true;                    // :27
 public:
  VariableVersion() : counter_(std::make_shared<VersionCounter>()) {} // :30 eager alloc
  explicit VariableVersion(bool enabled)   // :32 inference path
    : counter_(enabled? std::make_shared<VersionCounter>():nullptr), enabled_(enabled) {}
  uint32_t current_version() const { return counter_? counter_->current_version():0; } // :37
  bool is_enabled() const { return enabled_; }
  void bump(){ if(!enabled_) return; counter_->bump(); } // :44
  void reset(){ if(counter_) counter_->version_.store(0, std::memory_order_relaxed); } // :51
  bool shares_counter(const VariableVersion& o) const { return counter_ && counter_==o.counter_; } // :56
};
\end{lstlisting}

In \texttt{TensorImpl}:

\begin{lstlisting}[caption={and version wiring.}, label=lst:versionwire]
VariableVersion version_counter_{!InferenceMode::is_enabled()}; // :34 eager; disabled for inference tensors
bool inference_tensor_ = InferenceMode::is_enabled();           // :35
uint32_t version() const { // :265
  if(!version_counter_.is_enabled()) throw std::runtime_error("Inference tensors do not track version counter.");
  return version_counter_.current_version();
}
void bump_version(){ // :271
  if(inference_tensor_ && !InferenceMode::is_enabled())
    throw std::runtime_error("In-place update to an inference tensor outside inference mode is not allowed.");
  version_counter_.bump();
}
void share_version_counter(const TensorImpl& other){ // :289 alias for views
  inference_tensor_ = other.inference_tensor_;
  version_counter_ = other.inference_tensor_ ? VariableVersion(false) : other.version_counter_;
}
\end{lstlisting}

为什么急切：计数器在每个 \texttt{VariableVersion} 默认构造中分配，因此在任何突变之前创建的视图已经共享相同的 \texttt{VersionCounter} 对象作为其基础。延迟分配会让每个别名增长自己的计数器，默默地打破向后比较已保存版本与当前版本的 \texttt{SavedVariable} 防护。视图在创建时调用 \texttt{share\_version\_counter} ；克隆调用 \texttt{reset\_version} ，因此新存储从零开始，而源计数器继续增加。

\subsection{复制构造函数：什么是共享的，什么不是}
\label{sec:copyctor}

是明确的：

\begin{lstlisting}[caption={copy ctor sharing.}, label=lst:copyctor-mem]
TensorImpl::TensorImpl(const TensorImpl& other)
    : storage_offset_(other.storage_offset_),               // shared offset value (copied)
      sizes_and_strides_(other.sizes_and_strides_),         // copied inline/heap
      dtype_(other.dtype_), device_(other.device_),         // copied
      version_counter_(other.version_counter_),              // shared alias (same Counter*)
      inference_tensor_(other.inference_tensor_),            // copied
      is_contiguous_(other.is_contiguous_),                  // copied
      memory_format_(other.memory_format_),                  // copied
      storage_(other.storage_),                              // shared StorageImpl*
      shared_state_(other.shared_state_),                    // shared OneDNN cache
      sparse_state_(other.sparse_state_),                    // shared sparse components
      transform_value_(other.transform_value_),              // shared batched wrapper
      transform_batch_dim_(other.transform_batch_dim_),      // copied
      transform_level_(other.transform_level_),              // copied
      quantizer_(other.quantizer_) {}                        // shared immutable quantizer
// Note: autograd_meta_ intentionally absent -> copies start without grad history
\end{lstlisting}

共享内容： \texttt{storage\_} 因此副本是别名，直到写入触发调用者中的写入时复制； \texttt{shared\_state\_} 因此 OneDNN 描述符被重用； \texttt{sparse\_state\_} 因此 COO/CSR 分量张量不会重复； \texttt{transform\_value\_} 和 \texttt{quantizer\_} (, \texttt{:92}) 因此批处理和量化视图保持一致。不共享的内容：初始化列表中未提及 \texttt{autograd\_meta\_}，因此默认初始化为 \texttt{nullptr}——副本在没有等级历史记录的情况下重新开始，与中的显式注释相匹配。相同的规则适用于 \texttt{copy\_metadata\_from}，它在覆盖存储、大小和布局时保留目标的 \texttt{autograd\_meta\_}。

\begin{tpinsight}[Insight: Views Share Storage, History Does Not]
代码库中最重要的一个遗漏是 \texttt{TensorImpl} 复制构造函数中缺少 \texttt{autograd\_meta\_}。所有其他成员都在副本上共享——存储字节、版本计数器、稀疏组件、批处理包装器——因为这些描述了 \emph{what the data is}。历史不共享，因为它描述了 \emph{how this particular value was produced}：如果 \texttt{b = a} （副本，而不是视图）默默地继承了 \texttt{a} 的图，那么调用 \texttt{b.backward} 会将梯度流入 \texttt{a} 的父级，就好像 \texttt{a} 本身已被消耗一样——这是一个看不见的正确性错误。区别在于为什么 \texttt{detach} 被实现为复制加空而不是新的视图类型：共享 \texttt{TensorImpl} 会共享它不应该共享的东西。
\end{tpinsight}

\subsection{张量句柄：\texttt{shared\_ptr<TensorImpl>} 加上 479 个生成的方法}
\label{sec:tensorhandle}

是用户可见的句柄：

\begin{lstlisting}[caption={handle.}, label=lst:tensorhandle]
class P10_API Tensor {
 private:
  std::shared_ptr<TensorImpl> impl_;               // :82 the only data member
 public:
  Tensor() = default;
  explicit Tensor(std::shared_ptr<TensorImpl> impl) : impl_(std::move(impl)) {} // :88
  Tensor(const std::vector<int64_t>& sizes, DType dtype, const Device& device=Device()); // :90 allocate
  Tensor(Storage storage, const std::vector<int64_t>& sizes, DType dtype);               // :93 wrap
  // accessors: dim(), numel(), shape(), strides(), dtype(), device(), is_contiguous() ...
  // view/copy: view(), view_dtype(), as_strided(), select(), slice(), clone(), contiguous()
  // factory:   static empty/full/zeros/ones/arange/...  (generated)
  // generated surface:
#include "tensorplay/ops/TensorGenerated.h"        // :329  479 methods
};
\end{lstlisting}

include at 将代码生成的表面混合到类定义中。文件 \texttt{TensorGenerated.h} （6523 行，由生成）声明了来自（9\,415 行）的 479 个模式。每个条目都是一个成员声明，例如 \texttt{Tensor add(const Tensor\& other, Scalar alpha=1) const;} 或 \texttt{static Tensor zeros(...);}，通过 \texttt{Dispatcher} 和 \S\ref{sec:redispatch} 中记录的四阶段重新调度链（Vmap $\rightarrow$ Autocast $\rightarrow$ Autograd $\rightarrow$ Backend）进行调度。手写的核心保持很小（构造函数、数据访问器、视图原语、稀疏助手），而生成的标头提供了大部分操作表面。

辅助命名空间 \texttt{detail} 保存调度程序共享的低级克隆和连续原语：

\begin{lstlisting}[caption={detail clone surface.}, label=lst:detailclone]
namespace detail {
P10_API Tensor clone_impl(const Tensor& self, std::optional<MemoryFormat> memory_format = std::nullopt); // :339
P10_API Tensor contiguous_impl(const Tensor& self, int64_t memory_format);                               // :341
P10_API Tensor contiguous_clone(const Tensor& self);                                                       // :345 always contiguous
}
\end{lstlisting}

这些是自由函数，因此内核永远不会重新进入调度程序以获取纯字节副本；生成的 \texttt{Tensor::clone} 成员只是通过调度程序句柄转发到 \texttt{detail::clone\_impl} 。

\subsection{视图是零拷贝的：\texttt{as\_strided} 是单一瓶颈}
\label{sec:views}

每个别名操作---\texttt{view}、\texttt{transpose}、\texttt{select}、\texttt{slice}、\texttt{permute}、\texttt{view\_dtype}---通过\texttt{Tensor::as\_strided}进行传输：

\begin{lstlisting}[caption={\texttt{as\_strided}: reuse storage, share counter.}, label=lst:asstrided]
Tensor Tensor::as_strided(const std::vector<int64_t>& size,
                          const std::vector<int64_t>& stride,
                          std::optional<int64_t> storage_offset) const {
  if(!impl_) TP_THROW(RuntimeError, "Tensor not defined");
  if(size.size()!=stride.size()) TP_THROW(ValueError,"as_strided(): sizes and strides must have same length");
  for(int64_t v: size) if(v<0) TP_THROW(ValueError,"as_strided(): sizes must be non-negative");
  const int64_t offset = storage_offset.value_or(static_cast<int64_t>(impl_->storage_offset())); // :658 keep base offset when absent
  if(offset<0) TP_THROW(ValueError,"as_strided(): storage_offset must be non-negative");
  if(impl_->device().is_vulkan()){ // :666 Vulkan has no stride addressing -> materialize via kernel
    static const OperatorHandle handle = Dispatcher::singleton().findHandle("as_strided");
    return DispatchStub<Tensor,const Tensor&,const std::vector<int64_t>&,const std::vector<int64_t>&,std::optional<int64_t>>::call(handle, DispatchKey::Vulkan, *this, size, stride, std::optional<int64_t>(offset));
  }
  Tensor out = Tensor(impl_->storage(), size, stride, impl_->dtype(), static_cast<size_t>(offset)); // :675 reuse same Storage
  out.unsafeGetTensorImpl()->share_version_counter(*impl_); // :677 alias counter
  if(impl_->has_quantizer()) out.unsafeGetTensorImpl()->set_quantizer(impl_->quantizer()); // :680 quantizer rides along
  return out;
}
\end{lstlisting}

三个构建后步骤是视图契约：重用基础的 \texttt{Storage}，为基础的 \texttt{VersionCounter} 起别名，并且如果存在，则携带 \texttt{Quantizer}，以便视图在相同的代码上保持量化张量。这里使用的 \texttt{TensorImpl(Storage,...)} 构造函数通过 \texttt{SizesAndStrides::is\_contiguous} 根据提供的步长设置 \texttt{is\_contiguous\_}，因此即使对于别名张量，\texttt{is\_contiguous} 也保持正确。

来电者说明包装纸有多薄：

\begin{itemize}[leftmargin=1.2em]
\item \texttt{view} 通过 \texttt{SizesAndStrides::infer\_size} 推断 \texttt{-1} 变暗，通过 \texttt{compute\_view\_strides} 计算兼容的步幅，然后调用 \texttt{as\_strided(inferred,*stride)}。
\item \texttt{select} 删除一维，计算 \texttt{new\_offset = storage\_offset + index*stride[dim]}，从 \texttt{new\_strides} 中删除该维度，然后调用 \texttt{as\_strided}。
\item \texttt{slice} 固定 \texttt{step}，计算 \texttt{step}，将切片步长缩放 \texttt{step}，将偏移量提前 \texttt{start*stride[dim]}，然后调用 \texttt{as\_strided}。
\item 当元素大小匹配时 \texttt{view\_dtype} 保持存储和跨度；当它们不同时，它会重新调整 \texttt{new\_sizes.back}、领先步幅以及按比例 \texttt{dst\_esize} 偏移的元素。
\end{itemize}

这些路径中的任何一个都不会发生分配。 Vulkan 提前退出是唯一实现的分支：因为 Vulkan 有效负载没有跨步寻址，所以视图被降低到内核而不是别名存储。

\subsection{克隆和连续：分配新存储并重置版本}
\label{sec:clone}

(\texttt{detail::clone\_impl}) 是 \texttt{as\_strided} 的分配对应项：

\begin{lstlisting}[caption={\texttt{clone\_impl}: fresh storage, reset version.}, label=lst:cloneimpl]
Tensor detail::clone_impl(const Tensor& self, std::optional<MemoryFormat> memory_format){
  if(!self.defined()) return Tensor();
  if(self.is_sparse()){ // :891 sparse clones component-wise
    if(memory_format.has_value()) TP_THROW(RuntimeError,"unsupported memory format option ",toString(*memory_format));
    return Tensor::make_sparse_coo_tensor(self._indices().clone(), self._values().clone(),
                                          static_cast<std::vector<int64_t>>(self.shape()), self.is_coalesced());
  }
  const auto format = memory_format.value_or(MemoryFormat::Preserve); // :900
  const auto sizes_v = static_cast<std::vector<int64_t>>(self.shape());
  std::vector<int64_t> strides; MemoryFormat result_format = Contiguous; // :902
  if(format==Preserve){ // :904 preserve non-overlapping dense layout (e.g. transpose)
    const auto strides_v = static_cast<std::vector<int64_t>>(self.strides());
    if(SizesAndStrides::is_non_overlapping_and_dense(sizes_v, strides_v)) strides=strides_v;
    else strides=SizesAndStrides::compute_contiguous_strides(sizes_v);
  } else if(format==ChannelsLast){ strides=get_channels_last_strides(sizes_v); result_format=ChannelsLast; }
  else if(format==ChannelsLast3d){ strides=get_channels_last_strides(sizes_v); result_format=ChannelsLast3d; }
  else strides=SizesAndStrides::compute_contiguous_strides(sizes_v);
  Storage storage(static_cast<size_t>(self.numel())*self.itemsize(), getAllocator(self.device().type()), self.device()); // :924 fresh bytes
  auto out_impl = std::make_shared<TensorImpl>(storage, sizes_v, strides, self.dtype(), 0); // :926
  if(format==Preserve && strides==get_channels_last_strides(sizes_v) && strides!=SizesAndStrides::compute_contiguous_strides(sizes_v))
    result_format = sizes_v.size()==5? ChannelsLast3d : ChannelsLast; // :930 tag channels-last when preserved
  out_impl->set_memory_format(result_format);
  if(self.unsafeGetTensorImpl()->has_quantizer()) out_impl->set_quantizer(self.unsafeGetTensorImpl()->quantizer()); // :939
  Tensor t(std::move(out_impl));
  const auto rm_strides = SizesAndStrides::compute_contiguous_strides(sizes_v);
  const bool row_major = strides==rm_strides && static_cast<std::vector<int64_t>>(self.strides())==rm_strides; // :950
  if(self.device().is_cpu() && row_major) std::memcpy(t.data_ptr(), self.data_ptr(), size_t(self.numel())*self.itemsize()); // :953 fast path
  else t.copy_(self); // :957 layout-aware copy (handles channels-last etc.)
  t.unsafeGetTensorImpl()->reset_version(); // :961 clone starts unmutated (undo copy_ bump)
  return t;
}
\end{lstlisting}

关键决策：

\begin{itemize}[leftmargin=1.2em]
\item \textbf{Preserve vs.\ explicit format.} 使用 no/Preserve 格式，当输入不重叠且密集时（\texttt{SizesAndStrides::is\_non\_overlapping\_and\_dense} at），结果会保持输入的精确步幅（\texttt{SizesAndStrides::is\_non\_overlapping\_and\_dense} at），因此转置张量会克隆为转置张量，否则会回退到行优先（扩展/重叠输入）。显式 \texttt{ChannelsLast}/\texttt{ChannelsLast3d} 需要等级 4 或 5 并使用 \texttt{get\_channels\_last\_strides}。
\item \textbf{Fresh allocation.} \texttt{Storage(size\_t(self.numel)*itemsize, getAllocator(self.device.type), self.device)} 准确分配密集字节范围；稀疏克隆而是从克隆组件构建新的稀疏张量。
\item \textbf{Byte copy with fast path.} 当源和目标都是行主 CPU 张量时，克隆使用 \texttt{std::memcpy}；否则它使用 \texttt{copy\_}，它尊重通道最后一步和跨设备复制。 \texttt{contiguous\_clone} 是强制 \texttt{MemoryFormat::Contiguous} 的无条件变体，由需要平面指针算术的内核内部使用。
\item \textbf{Version reset.} 内部 \texttt{copy\_} 碰撞目的地计数器；克隆是一个新值，必须从零开始，因此 \texttt{reset\_version} 将其清除。源计数器未受影响。
\end{itemize}

\texttt{contiguous\_impl} 共享相同的结构：当在请求的格式中已经连续时，它返回 \texttt{self} 不变，否则以规范步长 \texttt{get\_strides\_for(sizes,format)} 实现一个新的连续张量，复制并重置版本。

\subsection{MemoryFormat：连续、ChannelsLast、ChannelsLast3d 和连续性测试}
\label{sec:memoryformat}

是布局枚举：

\begin{lstlisting}[caption={four tags, one sentinel.}, label=lst:memformat]
enum class MemoryFormat : int8_t { Contiguous=0, Preserve=1, ChannelsLast=2, ChannelsLast3d=3 };
// get_channels_last_strides: for [N,C,H,W] -> [C*H*W, 1, C*W, C]    // :26
//                   for [N,C,D,H,W] -> [C*D*H*W, 1, C*H*W, C*W, C]
inline std::vector<int64_t> get_channels_last_strides(const std::vector<int64_t>& sizes){ // :26
  const int64_t ndim=sizes.size(); std::vector<int64_t> strides(ndim);
  if(ndim<3){ int64_t prod=1; for(int64_t i=ndim-1;i>=0;--i){ strides[i]=prod; prod*=sizes[i]; } return strides; }
  strides[1]=1; strides[ndim-1]=sizes[1]; // :38 channel dim has stride 1
  for(int64_t i=ndim-2;i>=2;--i) strides[i]=strides[i+1]*sizes[i+1]; // :39 spatial dims
  strides[0]=strides[2]*sizes[2]; return strides; // :42 batch stride
}
inline std::vector<int64_t> get_strides_for(const std::vector<int64_t>& sizes, MemoryFormat f){ // :49
  switch(f){ case Contiguous: /* row-major */ case ChannelsLast: case ChannelsLast3d: return get_channels_last_strides(sizes); default: return {}; }
}
\end{lstlisting}

\texttt{Preserve} 是请求级哨兵，永远不会存储在活动的 \texttt{TensorImpl}----\texttt{memory\_format} 上，将其映射到 \texttt{Contiguous}，并在存储之前将其规范化为 \texttt{Contiguous}。阻塞格式的存储是 oneDNN 的工作，而不是步幅标签的工作：标签只是记录哪个规范步幅设置了张量声明。

\texttt{TensorImpl::is\_contiguous\_in} 是对布局敏感的内核很重要的谓词：

\begin{lstlisting}[caption={contiguity-in-format via exact stride compare.}, label=lst:iscontigin]
bool is_contiguous_in(MemoryFormat format) const {
  const auto sz = sizes_and_strides_.sizes().vec();
  const auto st = sizes_and_strides_.strides().vec();
  switch(format){
    case MemoryFormat::Preserve: return is_contiguous_;
    case MemoryFormat::ChannelsLast:
    case MemoryFormat::ChannelsLast3d: return st == get_channels_last_strides(sz);
    default: return is_contiguous_;
  }
}
\end{lstlisting}

排名规则在 \texttt{Tensor::is\_contiguous(MemoryFormat)} 中强制执行上一级：普通连续始终有意义，\texttt{ChannelsLast} 仅在排名~4 时有意义，\texttt{ChannelsLast3d} 在排名~5 时有意义，并且 \texttt{Preserve} 返回 true（它不施加布局约束）。变异路径是 \texttt{TensorImpl::set\_sizes\_and\_strides}，当提供的步幅与 \texttt{get\_channels\_last\_strides} 完全匹配时，它会从 \texttt{compute\_contiguous\_strides} 重新计算 \texttt{is\_contiguous\_}，并将 \texttt{memory\_format\_} 提升到 \texttt{ChannelsLast}/\texttt{ChannelsLast3d}。

\subsection{TensorImpl 上的 Autograd、Sparse、Vmap 和 Quant 扩展}
\label{sec:extensions}

Table~\ref{tab:tensorimplfields} 中的四个尾随组是扩展点，它们证明了平面结构而不是深层层次结构的合理性：

\begin{itemize}[leftmargin=1.2em]
\item \textbf{Autograd}。 \texttt{autograd\_meta\_:shared\_ptr<AutogradMetaBase>} 是每个张量的梯度钩子； \texttt{set\_autograd\_meta}/\texttt{autograd\_meta}/\texttt{has\_autograd\_meta} 公开它，\texttt{set\_requires\_grad} 拒绝推理张量。至关重要的是，每个复制操作都会删除它（请参阅 \S\ref{sec:copyctor}），因此复制开始时没有历史记录，并且只有 \texttt{autograd::make\_variable} 重新附加它。
\item \textbf{View flag}。 \texttt{is\_view\_} 标记张量是由视图操作创建的； \texttt{detach\_} 参考它来拒绝视图的突变。由 \texttt{as\_strided} 创建的视图不会在基础层中设置此标志——它是在实现公共视图时由更高级别的视图管道设置的——因此别名本身保留在 \texttt{as\_strided} 中，而标识位是单独管理的。
\item \textbf{OneDNN shared state}。两个不透明的 \texttt{shared\_ptr<void>} 保存一个 oneDNN 内存描述符和一个重新排序的内存缓存； \texttt{has\_onednn\_md} / \texttt{has\_onednn\_memory\_cache} 门他们的使用。共享是浅层的（指针复制），因此密集张量的副本重用相同的描述符。
\item \textbf{Sparse}。 \texttt{SparseState} 存储 \texttt{layout} (\texttt{0=COO, 1=CSR})、COO 的 \texttt{values} 和 CSR 的 \texttt{values}，以及 \texttt{sparse\_sizes}（逻辑大小）和 \texttt{coalesced}。安装稀疏状态时，密集存储会被清除，因此稀疏张量的逻辑形状不携带密集斑点。稀疏构造仅分配逻辑元数据，避免了大型嵌入的瞬态密集分配。
\item \textbf{Vmap wrapper}。 \texttt{transform\_value\_} 是物理张量（有批量维度），而 \texttt{sizes\_and\_strides\_} 是逻辑形状（没有批量维度）； \texttt{transform\_batch\_dim\_}/\texttt{transform\_level\_} 记录批次所在的位置以及哪个嵌套级别拥有包装器。 \texttt{is\_batched} 恰好是 \texttt{transform\_value\_!=nullptr}。
\item \textbf{Quantizer}。 \texttt{quantizer\_:shared\_ptr<Quantizer>} 是一个不可变的、共享的仿射映射（每个张量或每个通道尺度/零点，请参阅）。视图和克隆都传播它：视图通过 \texttt{as\_strided}，克隆通过 \texttt{clone\_impl}，在相同的代码上保持量化。
\end{itemize}

\subsection{放在一起：分配、视图和克隆演练}
\label{sec:walktogether}

考虑 CPU 上的 \texttt{a = tensorplay.ones([2,3])} 后跟 \texttt{b = a.view([3,2])} 和 \texttt{c = a.clone}：

\begin{enumerate}[leftmargin=1.2em]
\item \texttt{Tensor(\{2,3\}, Float32)} 构建一个 \texttt{TensorImpl(\{2,3\}, Float32, CPU)}，它分配 \texttt{Storage(2*3*4, getAllocator(CPU), CPU)} $\rightarrow$ \texttt{StorageImpl(24, CachingAllocator::instance)} $\rightarrow$ \texttt{CachingAllocator::allocate(24)} 舍入到 \texttt{(24+64)->64 = 128} 总字节，返回 \texttt{DataPtr(ptr+64, deleter, CPU)}。 \texttt{SizesAndStrides(\{2,3\})} 通过 \texttt{compute\_contiguous\_strides} 设置 \texttt{inline\_sizes=[2,3]} 和 \texttt{inline\_strides=[3,1]}。
\item \texttt{b = a.view([3,2])} 推断 \texttt{[3,2]}，计算 \texttt{compute\_view\_strides([2,3],[3,1],[3,2]) = [2,1]}，并调用 \texttt{as\_strided([3,2],[2,1])}，后者构造 \texttt{Tensor(a.storage, [3,2], [2,1], Float32, offset=0)}，重用相同的 \texttt{Storage} 句柄并为计数器设置别名。没有发生分配器调用； \texttt{a.storage.is\_same(b.storage)} 成立，\texttt{b.data\_ptr==a.data\_ptr + 0*itemsize} 成立。
\item \texttt{c = a.clone}（带有 \texttt{Preserve}）找到步长 \texttt{[3,1]} 非重叠密集，在同一设备上分配新的 \texttt{Storage(24, CachingAllocator::instance)}，在步长 \texttt{[3,1]} 处构建新的 \texttt{TensorImpl}，并通过 \texttt{memcpy} 复制 24 个字节，因为双方都是行优先。目标版本重置为零，而 \texttt{a} 的计数器不变。
\end{enumerate}

等级~4 的通道最后一个示例采用了不同的步幅路径：\texttt{x = tensorplay.randn([2,3,4,4])} 和 \texttt{x.contiguous(MemoryFormat.ChannelsLast)} 到达 \texttt{contiguous\_impl}，它在 \texttt{get\_channels\_last\_strides([2,3,4,4])=[48,1,12,3]} 处分配并通过 \texttt{copy\_} 进行复制，通过 \texttt{set\_memory\_format(ChannelsLast)} 保留标记，因此 \texttt{x.is\_contiguous(ChannelsLast)} 变为 true，\texttt{is\_contiguous} 为 false。


